\subsection{四元数代数是体的判别法}

设 \(\alpha,\beta,\gamma\) 是体 K 的元素，F 是型为 \((\alpha,\beta,\gamma)\) 的四元数代数。像 No.~1 中那样，我们把 F 的典范基记作 \((1,i,j,k)\)，并把 F 的子代数 \(K+Ki\) 记作 E。

命题 5. —— 下列性质等价：

\begin{enumerate}
\def\labelenumi{(\roman{enumi})}
\item
  四元数代数 \(\mathbf{F}\) 是一个体。
\item
  不存在非零元素 \(q \in \mathrm{F}\) 使得 \(\mathrm{T}_{\mathrm{F}}(q) = \mathrm{N}_{\mathrm{F}}(q) = 0\)。
\item
  F 的每个平方为零的元素都为零。
\item
  不存在任何非零向量 \((x,y,z,t)\)（在 \(K^{4}\) 中）使得
\end{enumerate}

\[
x ^ {2} + \beta x y - \alpha y ^ {2} - \gamma (z ^ {2} + \beta z t - \alpha t ^ {2}) = 0.
\]

\begin{enumerate}
\def\labelenumi{(\alph{enumi})}
\setcounter{enumi}{21}
\tightlist
\item
  不存在任何非零向量 \((x,y,z)\)（在 \(K^{3}\) 中）使得
\end{enumerate}

\[
x ^ {2} + \beta x y - \alpha y ^ {2} - \gamma z ^ {2} = 0.
\]

\begin{enumerate}
\def\labelenumi{(\roman{enumi})}
\setcounter{enumi}{5}
\tightlist
\item
  二次代数 E 是一个体，且 \(\gamma\) 不是 E 的任何元素的范数。
\end{enumerate}

F 的元素 q 是可逆的，当且仅当 \(\mathrm{N}_{\mathrm{F}}(q)\) 不为 0。(i) 与 (iv) 的等价性由 III, §2, No.~5, 第 445 页的公式 (31) 得出；(i) 蕴含 (iii) 且 (iv) 蕴含 (v) 是显然的。

(ii) 与 (iii) 的等价性由 VIII, 第 362 页，注 3 得出。

设 F 不是一个体。若 \(\gamma(4\alpha+\beta^{2})\neq0\)，则代数 F 是 K 上次数为 4 的中心单代数；它同构于代数 \(\mathbf{M}_{2}(\mathrm{K})\)（VIII, 第 120 页，定理 1），因而包含一个非零的平方为零的元素。若 \(\gamma=0\)，则我们有 \(j^{2}=0\)。若 \(4\alpha+\beta^{2}=0\)，则我们有 \((2i-\beta)^{2}=0\)，且当 K 的特征异于 2 时 \(2i-\beta\neq0\)。最后，若 K 的特征为 2 且 \(\beta\) 为零，则我们有 \(\mathrm{T}_{\mathrm{F}}(q)=0\)（对每个 \(q\in F\)；III, §2, No.~5, 第 445 页，公式 (31)）。由于 F 不是一个体，存在 F 的一个非零元素 q 使得 \(\mathrm{N}_{\mathrm{F}}(q)=0\)；我们有 \(q^{2}=0\)。这就证明了蕴含关系 \((\mathrm{iii})\Rightarrow(\mathrm{i})\)。

设 \(q = x + yi\) 是 E 的一个元素。我们有 \(\mathrm{N}_{\mathrm{E / K}}(q) = x^2 + \beta xy - \alpha y^2\)。设性质 (v) 成立。我们有 \(\mathrm{N}_{\mathrm{E / K}}(q) - \gamma \neq 0\)，且 \(\mathrm{N}_{\mathrm{E / K}}(q) \neq 0\)（若 \(q \neq 0\)），由此得到 (vi)。

最后，设性质 (vi) 成立。设 q 是 F 的一个非零元素；我们把它写成 \(u + vj\)（其中 u 与 v 在 E 中）。若 v 为零，则 q 是可逆的。若 v 不为零，则由 III, §2, No.~5, 第 443 页，公式 (24)，我们有 \(\mathrm{N}_{\mathrm{F}}(q) = \mathrm{N}_{\mathrm{F}}(v)\mathrm{N}_{\mathrm{F}}(v^{-1}u + j) = \mathrm{N}_{\mathrm{E}/\mathrm{K}}(v)\big(\mathrm{N}_{\mathrm{E}/\mathrm{K}}(v^{-1}u) - \gamma\big)\)。由于 \(\gamma\) 不是一个范数，\(\mathrm{N}_{\mathrm{F}}(q)\) 不为零，故 q 是可逆的。

注. —— 设四元数代数 F 是一个体。由等式 \(j^{2} = \gamma\) 推出，我们有 \(\gamma \neq 0\)。由 VIII, 第 363 页，命题 2，除非 K 的特征为 2 且 \(\beta\) 为零（此时代数 F 是交换的），F 的中心等于 K。

